\documentstyle{article}
\begin{document}
In section 14.3.6, 
 we say that expression (14.28) for the within cluster sum of squares
\begin{eqnarray*}
W(C)=\frac{1}{2}\sum_{k=1}^{K}\sum_{C(i)=k}\sum_{ C(i')=k} d(x_{i},x_{i'}).
\end{eqnarray*}  

can be written as expression (14.31):

\begin{eqnarray*}
W(C)  &=&\frac{1}{2}\sum_{k=1}^{K}\sum_{C(i)=k}\sum_{ C(i')=k}||x_i-x_{i'}||^2\cr
&=&\sum_{k=1}^{K}N_k\sum_{C(i)=k}||x_i-\bar x_k||^2.
\end{eqnarray*}
This is correct, but the latter is not the criterion that is minimized
by the usual
form of K-means clustering (Algorithm 14.1). Instead,
 K-means clustering minimizes the criterion
\begin{eqnarray}
W(C) =
\sum_{k=1}^{K}\sum_{C(i)=k}||x_i-\bar x_k||^2.
\label{one}
\end{eqnarray}
That is, there is no cluster size weight $N_k$,
where $N_k$ is the number of observations assigned to cluster $k$.
Intuitively, this occurs because (14.28) effectively gives a  weight
proportional to $N_k^2$ to  a cluster with $N_k$ members.
One could define a  modified version of (14.28)  that is  equivalent to
expression (\ref{one}) by including weights $1/N_k$.

Now in practice one could instead minimize (14.31): this would change
the assignment step (2) in Algorithm 14.1, which would no longer
simply assign a point to the cluster with the nearest centroid.
 Qualitatively, this would tend to
produce clusters with more equal cluster sizes than the standard
algorithm. This could in fact be a useful variation
on K-means clustering.

Thanks to Daniela Witten for finding this error.


\end{document}


